\scanpage{392}
\mspage{497}

it being understood, of course, that we let
$\mathcal G_{\rho,\sigma}^{*}$ operate on $G$ by inner automorphisms via the
factor $B'_0$. We must check that
\[
i_{\deleted{\unsure{B}}}(L_0)
=\{i_B(\lambda)i_G(\lambda^{-1})\mid\lambda\in\Gamma L_0\}
\]
is invariant under
$\deleted{\unsure{\Gamma T_\rho},\unsure{\Gamma T_\sigma}\unsure{i_T}}$
\inserted{\unsure{the}}
$i_\rho(x)$, $i_\sigma(x)$, $i_B(x)$ and $i_G(x)$.
--- For the first two, this is trivial --- with trivial action on $L_0$.

If $x\in\Gamma B'_0$, we will have
\[
\begin{aligned}
i_B(x)(i(\lambda))
&=i_B(x)(i_B(\lambda)i_G(\lambda^{-1}))=\\
&=\underbrace{i_B(x)(i_B(\lambda))}
_{i_B(\deleted{\unsure{\lambda|x}}x(\lambda))}
\cdot
\underbrace{i_B(x)(i_G(\lambda^{-1}))}
_{\substack{
i_G(x)(i_G(\lambda^{-1}))\\
=i_G(x(\lambda)^{-1}]
}}\\
&=i(x(\lambda))
\end{aligned}
\]
\margin{We \unsure{see} at once that $i(\lambda)$
commutes with $G$, hence is invariant under $G$ !}
which shows that $\deleted{\unsure{\mathbb G_a|L_0}}$
\[
\diagram{\begin{tikzcd}[ampersand replacement=\&,column sep=large]
L_0\arrow[r,"i"] \&
\mathcal G_{\rho,\sigma}^{*}\cdot G
=\widetilde S\mathcal G_{\rho,\sigma}^{*}
\end{tikzcd}}
\]
commutes with the natural action of $B'_0$ on both sides
(on the second via $\unsure{\pi_B}$,
\[
\diagram{\begin{tikzcd}[ampersand replacement=\&,column sep=large]
i_B:B'_0\arrow[r] \&
\widetilde S\underset{\deleted{\unsure{0}}}{\mathcal M}_{\rho,\sigma}^{\unsure{\flat|*|\natural}}
\end{tikzcd}}
\]
). Finally
\[
i_G(\deleted{\unsure{f}}\inserted{\unsure{x}})(i(\lambda))
=i_G(x)\cdot(i_B(\lambda)i_G(\lambda^{-1}))\cdot i_G(x)^{-1}
\]
now by definition, $i_B(\lambda)$ and $i_G(\lambda)$ operate
in the same way on $G$, hence
$i(\lambda)=i_B(\lambda)i_G(\lambda^{-1})$
operates trivially on it, hence the last factor
$i_G(x)^{-1}$ can \unsure{jump} \unsure{over|under} the factor
$i(\lambda)$, we find
\[
i_G(x)i_G(x)^{-1}i(\lambda)=i(\lambda)
\]
\margin{O.K.}
\margin{391}

% Printed archive footer (not autograph).
\begin{center}
391\\
\texttt{https://grothendieck.umontpellier.fr}\\
Université de Montpellier, UMR CNRS 5149, Case courrier 051, Place Eugène Bataillon\\
34095 Montpellier cedex 5 - France
\end{center}

\scanpage{393}

\uline{Remark} $\widetilde S\mathcal M_{\rho,\sigma}$ is a flat group of relative dimension
$\underbrace{2}_{B'_0}+\underbrace{2}_{T_\rho}+\underbrace{2}_{N'_\sigma}+\underbrace{4}_{G}=10$,
$S\mathcal M_{\rho,\sigma}$ is of relative dimension
$\underbrace{\unsure{4}}_{S\mathcal G_{\rho,\sigma}^{*}}+\underbrace{3}_{SG}=7$,
on passing to the quotient by $\underset{\simeq\mathbb G_a}{i(L_0)}$
\unsure{gives|to give} groups \deleted{\unsure{this}} denoted
$\widetilde{\mathcal M}_{\rho,\sigma}$, $\mathcal M_{\rho,\sigma}$,
of relative dimension $9$ resp. $6$, giving rise to a diagram
\[
(43)\qquad
\diagram{\begin{tikzcd}[ampersand replacement=\&,column sep=large,row sep=large]
1 \arrow[r] \& \mathbb G_a \arrow[r,"\lambda\mapsto\varepsilon_0(\lambda)"] \arrow[d,equal]
  \& S\mathcal M_{\rho,\sigma} \arrow[r] \arrow[d,hook,"\unsure{\mathrm{can.}}"']
  \& \mathcal M_{\rho,\sigma} \arrow[r] \arrow[d,hook] \& 1 \\
1 \arrow[r] \& \mathbb G_a \arrow[r]
  \& \widetilde S\mathcal M_{\rho,\sigma} \arrow[r]
  \& \unsure{\mathcal M_{\rho,\sigma}|\widetilde{\mathcal M}_{\rho,\sigma}} \arrow[r,no head] \& 1
\end{tikzcd}}
\]

The passage to the quotient presents no difficulties --- the quotients are in fact
\unsure{representable} \deleted{\unsure{by the|by}}\inserted{as} $1/2$ direct products
\[
(44)\qquad
\left\{\begin{aligned}
\deleted{\unsure{S}}\mathcal M_{\rho,\sigma}
 &\simeq \deleted{\unsure{T_0}}\,S\mathcal G_{\rho,\sigma}^{*}(0)\cdot SG\\
\widetilde{\mathcal M}_{\rho,\sigma}^{\deleted{\unsure{\sim}}}
 &\simeq \mathcal G_{\rho,\sigma}^{*\deleted{\unsure{\flat|\natural|\circ}}}(\unsure{0|\omega})\cdot\deleted{\unsure{S}}G
\end{aligned}\right.
\]
where the subgroup schemes
$S\mathcal G_{\rho,\sigma}^{*\unsure{\flat|\natural|\circ}}(\unsure{0|\omega})$ and
$\mathcal G_{\rho,\sigma}^{*\unsure{\flat|\natural|\circ}}(\unsure{0|\omega})$ of
$S\mathcal G_{\rho,\sigma}^{*}$, $\mathcal G_{\rho,\sigma}^{*}$
are defined in $\boxed{\text{\uline{VIII}}}$,
$\mathcal G_{\rho,\sigma}^{*\unsure{\flat|\natural|\circ}}(\unsure{0|\omega})$ being identified with
$T_0\times T_\rho\times N'_\sigma$, and
$S\mathcal G_{\rho,\sigma}^{*\unsure{\flat|\natural|\circ}}(0)$ with its intersection
with $S\mathcal G_{\rho,\sigma}^{*}$ (defined by the supplementary conditions $(20)$ on page $487'$).

\margin{392}

% Printed archive footer (not autograph).
\begin{center}
392\\
\texttt{https://grothendieck.umontpellier.fr}\\
Université de Montpellier, UMR CNRS 5149, Case courrier 051, Place Eugène Bataillon\\
34095 Montpellier cedex 5 - France
\end{center}

\scanpage{394}
\mspage{498}

Returning now to the arithmetic situation, we find
\inserted{by passage to the quotient via $(37)$} a canonical \unsure{homomorphism}
\margin{in view of $(13)$}
\[
(45)\qquad
\diagram{\begin{tikzcd}[ampersand replacement=\&,column sep=large]
\deleted{\unsure{S_0}}\mathcal M_{0,3}^{!\sim} \arrow[r]
\& \underbrace{\mathcal M_{\rho,\sigma}(\widehat{\mathbb Z})}_{\begin{gathered}
\simeq\\
\underbrace{\deleted{\unsure{B'_0}}\unsure{T_0|\widehat T_0}\times
T_\rho(\widehat{\mathbb Z})\times T_\sigma(\widehat{\mathbb Z})}_{\begin{gathered}
\text{operating on}\\
\mathrm{Sl}(2,\widehat{\mathbb Z})\ \text{via the}\\
\text{factor }T_0
\end{gathered}}\cdot\mathrm{Sl}(2,\widehat{\mathbb Z})
\end{gathered}}/\Sigma
\end{tikzcd}}
\]
$\Sigma$ being the subgroup $\simeq\{\pm1\}^3$, product
of three factors $\mathbb Z/2\mathbb Z$ included resp.$^{\unsure{t}}$
in $T_\rho(\widehat{\mathbb Z})$, $T_\sigma(\widehat{\mathbb Z})$
and in $\mathrm{Sl}(2,\widehat{\mathbb Z})$.

\uline{Remark} This homomorphism was obtained here,
\inserted{essentially,} with the help of the homomorphism from
$\mathcal M_{0,3}^{!\sim}\unsure{(0)}$
\[
(46)\qquad
\diagram{\begin{tikzcd}[ampersand replacement=\&,column sep=large,row sep=huge]
\begin{gathered}
\mathcal M_{0,3}^{!\sim}(0)\\
\text{subgroup}\\
\text{of the automorphisms}\\
\text{``normalized''}
\end{gathered} \arrow[dr]
\& \begin{gathered}
=\operatorname{Norm}_{\mathcal M_{0,3}^{!\sim}}(L_0)\\
\text{(cf. }\S46,\ \text{\uline{V}}\text{)}
\end{gathered} \arrow[dr,-{Circle[length=2pt,width=2pt]}] \& \\
\& S\mathcal G_{\rho,\sigma}^{*\unsure{\flat|\natural|\circ}}(0)(\widehat{\mathbb Z})/\Sigma_0
  \arrow[r,hook,no head]
\& S\mathcal G_{\rho,\sigma}^{*}(\widehat{\mathbb Z})/\Sigma_0
\end{tikzcd}}
\]
and of the composite
\[
\diagram{\begin{tikzcd}[ampersand replacement=\&,column sep=large,row sep=small]
\deleted{\unsure{\mathcal G}}\widehat\Pi_{0,3}
  \arrow[r,hook,no head]
\& \widehat{\unsure{\mathcal G|\mathcal C}}_{0,3}^{+}
  \arrow[r] \arrow[d,phantom,"\unsure{\cup|\wr}" description]
\& \mathrm{Sl}(2,\widehat{\mathbb Z})' \\
\& \mathrm{Sl}(2,\mathbb Z)^{\prime\wedge} \&
\end{tikzcd}}
\]
by writing
\[
(47)\qquad
\mathcal M_{0,3}^{!\sim}\simeq\mathcal M_{0,3}^{!\sim}\cdot\widehat\Pi_{0,3}\quad\unsure{!|:}
\]

\margin{393}

% Printed archive footer (not autograph).
\begin{center}
393\\
\texttt{https://grothendieck.umontpellier.fr}\\
Université de Montpellier, UMR CNRS 5149, Case courrier 051, Place Eugène Bataillon\\
34095 Montpellier cedex 5 - France
\end{center}

\scanpage{395}

But we also have a semidirect product
\[
(48)\qquad
\mathcal M_{0,3}^{\sim}\simeq
\mathcal M_{0,3}^{!\sim}(\unsure{0|\omega})\cdot
\unsure{\mathcal G|\mathcal C}_{0,3}^{+}
\]
which, combined directly with
\[
\diagram{\begin{tikzcd}[ampersand replacement=\&,column sep=large]
\widehat{\unsure{\mathcal G|\mathcal C}}_{0,3}^{+}\arrow[r]
\&\mathrm{Sl}(2,\widehat{\mathbb Z})'
\end{tikzcd}}
\]
gives a homomorphism
\[
(49)\qquad
\diagram{\begin{tikzcd}[ampersand replacement=\&,column sep=large]
\mathcal M_{0,3}^{\sim}\arrow[r]
\& \bigl(\mathcal M_{\rho,\sigma}(\widehat{\mathbb Z})
\simeq S\mathcal G_{\rho,\sigma}^{*}(\widehat{\mathbb Z})\cdot
\mathrm{Sl}(2,\widehat{\mathbb Z})\bigr)/\unsure{\Sigma}
\end{tikzcd}}
\]
which extends the homomorphism $(45)$, defined first on its normal
subgroup $\mathcal M_{0,3}^{!\sim}$ of index $6$.

\margin{394}

% Printed archive footer (not autograph).
\begin{center}
394\\
\texttt{https://grothendieck.umontpellier.fr}\\
Université de Montpellier, UMR CNRS 5149, Case courrier 051, Place Eugène Bataillon\\
34095 Montpellier cedex 5 - France
\end{center}

\scanpage{396}
\mspage{499}

$\overline{\text{\uline{X}}}$\quad Study of the group
$\mathcal M_{\rho,\sigma}$ and of $\mathcal M_{\rho,\sigma}\to\mathrm{Gl}(2)$

By construction of the semidirect product
\[
(1)\qquad
\diagram{\begin{tikzcd}[ampersand replacement=\&,column sep=small,row sep=huge]
\mathcal M_{\rho,\sigma}^{\sim} \arrow[r,phantom,"=" description]
\& \mathcal G_{\rho,\sigma}^{*\unsure{\flat|\natural|\circ}}(\unsure{\omega|0})\cdot G
  \arrow[r,phantom,"\simeq" description] \arrow[d,phantom,"\cup" description]
\& \underbrace{(T_0\times T_\rho\times N'_\sigma)}_{\begin{gathered}
\text{operating on }G\\ \text{via }T_0\hookrightarrow G
\end{gathered}}\cdot G \\
\mathcal M_{\rho,\sigma} \arrow[r,phantom,"=" description]
\& S\mathcal G_{\rho,\sigma}^{*\unsure{\flat|\natural|\circ}}(\unsure{\omega|0})\cdot SG \arrow[r,hook]
\& \underbrace{(T_\rho\times N'_\sigma)}_{\begin{gathered}
\text{operating on }G\\
\text{via }N'_\sigma\xrightarrow{\det}\mathbb G_m
\xleftarrow[\sim]{\det}T_0\hookrightarrow G
\end{gathered}}\cdot G
\end{tikzcd}}
\]
we have a group homomorphism
\[
(2)\qquad
\left\{\begin{gathered}
\diagram{\begin{tikzcd}[ampersand replacement=\&,column sep=large]
\mathcal M_{\rho,\sigma}^{\sim}\arrow[r,"\Theta"] \& G
\end{tikzcd}}\\
\Theta\bigl(i_B(x)i_{\deleted{\unsure{\rho|B}}\inserted{\rho}}(y)
\unsure{i_\sigma|i'_\sigma}(z)i_G(f)\bigr)=x.f
\end{gathered}\right.
\]
\[
(x\in\Gamma T_0,\quad y\in\Gamma T_\rho,\quad
z\in\Gamma\unsure{N'_\sigma},\quad f\in\Gamma G).
\]

The composite homomorphism
\[
(3)\qquad
\diagram{\begin{tikzcd}[ampersand replacement=\&,column sep=large,row sep=large]
S\mathcal M_{\rho,\sigma}^{\sim}\arrow[r,no head]
  \arrow[d,phantom,"\unsure{\wr}" description]
\&\mathcal M_{\rho,\sigma}^{\sim}\arrow[r,"\Theta"] \& G\\
\underbrace{\deleted{\unsure{\mathcal M}}\mathcal G_{\rho,\sigma}^{*\unsure{\flat|\natural|\circ}}}_{B'_0\times T_\rho\times N'_\sigma}\cdot G
\&\&
\end{tikzcd}}
\]
\margin{(composite still denoted $\Theta$)}
\deleted{is} described by
\[
\Theta\bigl(i_B(x)i_\rho(y)i_\sigma(z)i_G(\unsure{f|\varphi})\bigr)=x.f
\]
and $\Theta$ will indeed be trivial on the
\[
i(\lambda)=i_B(\lambda)i_G(\lambda^{-1})\qquad(\lambda\in\Gamma L_0)
\]
and therefore passes to the quotient as $\Theta:\mathcal M_{\rho,\sigma}^{\sim}\to G$.

\margin{395}

% Printed archive footer (not autograph).
\begin{center}
395\\
\texttt{https://grothendieck.umontpellier.fr}\\
Université de Montpellier, UMR CNRS 5149, Case courrier 051, Place Eugène Bataillon\\
34095 Montpellier cedex 5 - France
\end{center}